\documentclass[10pt]{letter}

\usepackage[a4paper,margin=0.85in]{geometry}
\usepackage[T1]{fontenc}
\usepackage[utf8]{inputenc}
\usepackage{lmodern}
\usepackage{microtype}

\signature{Salah-Eddin Gherbi}
\address{
Salah-Eddin Gherbi\\
Independent Researcher\\
Trowbridge, United Kingdom\\
Email: salahealer@gmail.com\\
ORCID: 0009-0005-4017-1095
}

\begin{document}

\begin{letter}{
The Editor\\
International Journal of Applied and Computational Mathematics
}

\opening{Dear Editor,}

I am pleased to submit ``A Certified Two-Layer Presence Census of Rao's
Spherical \'{S}r\={\i} Yantra Constraint System: Multiplicity, Per-Root
Geometric Validity, and a Rule-Governed Discovery Program'' for consideration
in the \emph{International Journal of Applied and Computational Mathematics}.

C. S. Rao (1998) formalized the spherical \'{S}r\={\i} Yantra as a system of
nonlinear equations in six angular variables subject to twenty geometric
constraint conditions. This manuscript reports a certified computational
presence census of that system over the registered universe of 3044 well-posed
choices of six constraints. The census certifies 888 feasible subsystems
supporting 968 Krawczyk-validated roots, and each certified root is separately
classified by a registered Gate-4/Rao geometric-validity test. One consequence
is that geometric validity is an irreducibly per-root property: 26 subsystems
carry both valid and rejected certified roots.

The work is methodologically aligned with the scope of the
\emph{International Journal of Applied and Computational Mathematics}. Its core
is validated numerical computation: interval-arithmetic root certification, a
deterministic domain-wide discovery generator, a preregistered discovery
protocol, and a complete reproducibility trail in which the reported results
are regenerable from frozen, hash-identified code and publicly deposited data.
Importantly, unresolved subsystems are reported with scoped labels rather than
overclaimed as infeasible; the paper certifies presence, not absence. I believe
this combination of nonlinear geometric constraints, interval validation,
computational enumeration, and reproducible numerical methodology makes the
manuscript a suitable contribution to the journal.

For transparency, a companion manuscript by the same author, which establishes
the underlying plane--sphere reduction as a theorem, is currently under review
at another journal. A related manuscript on the plane form of the system is also
under review at the \emph{Indian Journal of History of Science}. These
manuscripts report distinct results with no duplication of text or findings:
the plane manuscript concerns the certified enumeration of the planar constraint
system, the companion manuscript concerns the validated plane--sphere bridge
theorem, and the present manuscript concerns the certified spherical presence
census. The present manuscript is self-contained and does not depend on the
others. I have attached the companion manuscript as a related file for
editorial reference.

All supporting materials, including the certified census dataset, signed
manifest and hash lineage, and the public atlas of certified figures, are
publicly deposited with DOIs prior to submission. The manuscript includes a
declaration on the use of large language models as computational,
organizational, and drafting aids. I take full responsibility for all
mathematical claims, computations, content, and conclusions.

The manuscript is not under consideration elsewhere. I confirm that I am the
sole author.

Thank you for considering this work.

\vspace{0.8em}

\noindent Yours sincerely,\\[1.5em]
Salah-Eddin Gherbi

\end{letter}

\end{document}